\documentclass[11pt]{article}
\usepackage[margin=1in]{geometry}
\usepackage{amsmath,amssymb}
\usepackage[expansion=false]{microtype}
\usepackage[hidelinks]{hyperref}
\usepackage{enumitem}
\usepackage{array}
\usepackage{xcolor}
\usepackage[compact]{titlesec}
\titleformat{\section}{\Large\bfseries}{\thesection}{0.7em}{}
\titlespacing*{\section}{0pt}{12pt}{4pt}
\setlist{nosep,leftmargin=1.4em}
\setlist[itemize]{label=$\bullet$}
\setlength{\parskip}{0.4em}
\setlength{\parindent}{0pt}
\makeatletter
\renewcommand\paragraph{\@startsection{paragraph}{4}{\z@}{0.7ex \@plus 0.3ex}{-1em}{\normalfont\normalsize\bfseries}}
\makeatother
\setlength{\abovedisplayskip}{6pt}\setlength{\belowdisplayskip}{6pt}

\newcommand{\E}{\mathbb{E}}
% Placeholder for content we still need to write; remove all before submitting.
\newcommand{\todo}[1]{\textcolor{red}{\textbf{[TODO:} #1\textbf{]}}}

\begin{document}

\begin{center}
{\LARGE Checkpoint 3: Adaptive Negotiation Training\\[2pt] Through Human-Centered RL}\\[8pt]
{\large Yash Ravipati, Mahdi Nassereddine, Vincent Song}\\[2pt]
CS2822r, Fall 2026
\end{center}
\vspace{-0.6em}

% ===========================================================================
\section{Research question}

Our question has narrowed since Checkpoint~2: does practicing an upcoming negotiation
against an LLM counterpart, with feedback after each round, help someone do better in
that negotiation? The user uploads the case, the tool finds where they are weak in it,
and it sets up practice rounds that target those moments. We focus on institutional
conflict and coordination, as in class Simulations~3 and~4.

Following the Checkpoint~2 feedback, we replaced our job-offer bargaining errors with
skills taught in class, such as building rapport and finding the interests behind a
position, and we no longer use LLM-simulated trainees as evidence that training works.
We ask:
\begin{itemize}
\item \textbf{Q1, diagnosis:} does the tool spot a user's weaknesses as well as a human
annotator?
\item \textbf{Q2, steering:} does focusing the counterpart on a skill create chances to
practice it, without breaking the case?
\item \textbf{Q3, preparation:} do people who prepare with the tool improve more on the
case than those who prepare without it?
\end{itemize}

\paragraph{Model:}
A case $c$ includes each side's private materials, the constraints and approvals a deal
needs, and a rubric of $J=6$ skills (Table~\ref{tab:rubric}). For each user message in a
practice round, we score every skill it gave a chance to use, $q\in[0,1]$, from $0$
(missed) to $1$ (effective). A round's score $Q_j(\tau_t)$ is the average over its
$n_{t,j}$ scored chances, or blank if there were none. Before each round, the coach picks
a difficulty $\mu_t\in\{1,\dots,5\}$ and a focus skill $f_t$, or none. It tracks the
user's hidden skill level $s_{t,j}$ with a Beta distribution, where a score $q$ counts as
$q$ successes and $1-q$ failures and older evidence fades by $\lambda\in(0,1]$:
\[
\alpha_{t+1,j}=\alpha_{0,j}+\lambda(\alpha_{t,j}-\alpha_{0,j})+v_{t,j},\qquad
\beta_{t+1,j}=\beta_{0,j}+\lambda(\beta_{t,j}-\beta_{0,j})+n_{t,j}-v_{t,j},
\]
with $v_{t,j}=\sum_w q_{w,j}$. The next round targets the skill that looks weakest,
favoring skills with little evidence. Success means case performance, the expected score
on $c$ at a fixed difficulty $\mu_0$ with no focus or feedback, $G_j(s;c)=\E[Q_j(\tau)]$;
we want it to rise more with the tool than without.

\begin{table}[t]
\centering
\footnotesize
\renewcommand{\arraystretch}{1.0}
\begin{tabular}{*{5}{>{\raggedright\arraybackslash}p{0.172\linewidth}}}
\hline
\textbf{Skill} & \textbf{Opportunity arises when} & \textbf{Missed (0)} &
\textbf{Partial ($\approx$0.5)} & \textbf{Effective (1)} \\
\hline
Rapport and shared values &
A concern, tension, or stated value gives a reason for constructive engagement. &
Dismisses the concern, or uses a mission statement as an unsupported accusation. &
Acknowledges the concern or value without connecting it to the discussion. &
Responds respectfully and connects the concern or value to a shared purpose. \\
Interests, constraints, and decision-making &
A demand, ambiguous restriction, or pending approval needs clarification. &
Treats a demand as a fixed constraint, or overlooks who must approve. &
Asks a generic question without probing the relevant motivation or approval. &
Explores and checks the underlying interest, constraint, or decision dependency. \\
Facts and shared criteria &
A claim is disputed, a fact is unknown, or a standard must be chosen. &
Treats an uncertain claim as established, or ignores relevant evidence. &
Notes the uncertainty or a criterion without resolving or applying it. &
Checks the evidence, or agrees how to apply a suitable criterion. \\
Leverage &
A request or obstacle makes the learner's contribution or alternative consequential. &
Misses supported leverage, or exaggerates a dependency or alternative. &
Recognizes the contribution but does not tie it to a proposal or request. &
Accurately connects the contribution or alternative to a relevant proposal. \\
Options for multiple interests &
Expressed interests or an impasse allow another arrangement to be explored. &
Repeats a position, or proposes an option that ignores known constraints. &
Suggests an alternative without linking it to interests or checking feasibility. &
Develops and tests a feasible option that addresses several expressed interests. \\
Commitments and follow-through &
A tentative understanding or promise needs clarifying before action or closure. &
Treats an unsupported assurance as settled, or leaves a key dependency unexamined. &
Records a next step but leaves the owner, approval, deadline, or contingency unclear. &
Clarifies who does what, by when, with which approvals, or records why approval is pending. \\
\hline
\end{tabular}
\caption{Scoring rubric for institutional conflict and coordination. Scores can be any
value from 0 to 1; the columns are reference points.}
\label{tab:rubric}
\end{table}

\paragraph{How a transcript is scored:}
After a round, the scorer reviews each user message seeing only the user's own materials
and the conversation so far, so each move is judged on what the user knew then. For each
skill, it notes whether there was a chance to use it, scores it with
Table~\ref{tab:rubric}, and quotes the exact words it relied on; code checks the quotes
are real. Skipping a skill counts as a miss only if the moment called for it. For
example, if in Simulation~4 the other side says a referral needs headquarters' approval,
``Then we will proceed without you'' scores near 0 on interests and decision-making,
while ``What would headquarters need to see to approve this?'' scores near 1. Claude does
the scoring in the tool; human annotators use the same rubric and format, so the two can
be compared directly.

% ===========================================================================
\section{What we have completed}

The app covers the whole preparation loop: upload the case, read a prep brief, practice,
get feedback, and see a skill profile that picks the next round. Since Checkpoint~2 we
have rebuilt it around the rubric.
\begin{itemize}
\item \textbf{Cases from uploads:} uploaded materials become a case with its negotiation
type, constraints, required approvals, and the six-skill rubric. Point tables are
optional. A conversion script sets up Simulations~3 and~4 this way.
\item \textbf{A steerable counterpart:} Claude plays the other side using only that
side's materials. Each skill has instructions for creating practice chances, such as
raising an unsettled approval to practice interests and decision-making. The counterpart
sticks to the case's facts, and being friendly never earns a free concession.
\item \textbf{Scoring and feedback:} the message-by-message scoring above, plus a report
organized by skill that shows what the user said, why it mattered, and a better option.
\item \textbf{Profile and next round:} a Beta estimate for each skill, shown with how
much evidence backs it, used to choose the next focus.
\item \textbf{Evaluation tools:} import and export of human annotations, and code that
compares the system with them, measures steering, and summarizes feedback ratings.
\end{itemize}
The counterpart never sees the user's confidential materials, and the user doesn't see
the counterpart's until the debrief. The system has 165 automated tests.

% ===========================================================================
\section{How the tool has done}

So far we have the stress tests from Checkpoint~2 and one real use in class. Our main
focus is the evaluations with human judges and real participants that follow, which have
not run yet.

\paragraph{Algorithmic stress testing:}
For Checkpoint~2 we ran 450 offline negotiations between our rule-based counterpart,
which follows the same concession rules as the LLM counterpart, and scripted users with
three styles. The counterpart never accepted a deal below its walk-away point, and harder
levels consistently cost users points (Spearman $\rho=-0.82$). These tests show the
mechanics work, not that the tool helps anyone prepare.

\paragraph{In-class Simulation 4:}
One of us used the tool to prepare for Simulation~4. The biggest problem was time: the
class gave us 15 minutes, which isn't enough to practice, get feedback, and practice
again, and that cycle is where the training happens. Even so, the tool helped in two
ways. First, it laid out the main issues, the options for each, and how many points each
option was worth, which made it clear what to prioritize. These numbers come from code
working on the user's own sheet, not from the LLM. Second, the built-in assistant was
very useful for quick questions while preparing. It answers only from the user's own
materials, so it can clear up facts without giving away the other side's information.
Working through the case this way also kept the details fresh going into the negotiation.
We haven't yet tried a longer preparation window; the study below will.

\paragraph{Weakness identification (Q1, planned):}
Two of us will separately label 20--30 transcripts from Simulations~3 and~4 in the
system's format, never labeling a transcript we played. We will measure how often the
system finds the same chances and weaknesses as we do (precision and recall) and how far
its scores are from ours. We will compare it with a keyword detector, with a
general-purpose model that doesn't use our rubric, and with how well our two labelers
agree, which is about the best we can expect. We will also remove one skill from single
messages and check that only that skill's score drops.

\paragraph{Steering (Q2, planned):}
We will take about ten conversations stopped partway through and generate the
counterpart's next message under each focus and with none. Raters who don't know the
focus will judge which practice chances each message creates and whether it stays true to
the case. Steering works if each focus creates more chances for its own skill than no
focus does. We will also check whether adversarial messages can make the counterpart leak
its confidential information.

\paragraph{Preparation (Q3, planned):}
Scores from the practice rounds can't show on their own that preparation worked: the
focus and difficulty change between rounds, people improve with any practice, and the
system would be grading itself. Our real test compares before and after on the same
case. Each participant first negotiates a new case at a fixed difficulty, with no focus
and no feedback. We then randomly split them into two groups with the same preparation
time. One group uses our tool. The other gets the same case files and a general-purpose
chatbot like ChatGPT. Then everyone negotiates the case again under the same conditions
and, when possible, against a real person. Claude scores every transcript with
Table~\ref{tab:rubric}. Because the tool trains people against that same scorer, a team
member or TA also scores a random 20--30\% blind to group and order, to check the gains
hold up for a human. We will also record the outcome (agreement, constraints respected,
clear commitments) and, after negotiations against people, how the other side rated trust
and fairness. With 6--10 participants, this will be early evidence, not proof.

% ===========================================================================
\section{What we keep and what we change}

\paragraph{Keep:}
Practicing the user's own upcoming case, a counterpart that only knows its own side's
materials, and feedback after the session rather than during it, so users build skill
instead of leaning on the tool. We also keep the Beta estimate for each skill, which now
takes partial scores instead of just right or wrong. Simulation~4 gave us reason to keep
the prep brief and the assistant, the most useful parts under time pressure.

\paragraph{Change:}
\begin{itemize}
\item \textbf{Skills from class:} the six-skill rubric replaces the bargaining errors.
\item \textbf{A case-specific goal:} we measure improvement on the case being prepared,
not negotiation skill in general.
\item \textbf{Evaluation with people:} human labels and a before-and-after study replace
LLM-simulated trainees.
\item \textbf{Short preparation windows:} when time is short, the tool should lead with
the prep brief and assistant. When there's more time, it should push users to start
early enough for several practice rounds.
\end{itemize}

\paragraph{Next steps:}
We will run the diagnosis and steering evaluations first, then the before-and-after
study. Alongside that, we will improve the interface, make preparation smoother when time
is short, and ask the instructor or a TA to review the rubric.

\end{document}
